\documentclass[12pt]{article}
\usepackage[utf8]{inputenc}
\usepackage[T1]{fontenc}
\usepackage{hyperref}
\title{Tilung Language Resources for K-12 Curriculum Development}
\author{Compiled by Claude Sonnet 5 (claude-sonnet-5) for Andr\'as Kornai}
\date{}
\begin{document}
\maketitle

\section{Introduction}
Tilung (ISO 639-3: tij) is a Kiranti language of Khotang district, Nepal. It is severely endangered: one source reports only seven elderly fluent speakers, Wikipedia gives 20, and younger people speak Nepali. Only academic materials were found.

\section{Primary Educational Resources}
None found.

\section{Linguistic and Reference Materials}
\begin{itemize}
\item Wikipedia entry: \url{https://en.wikipedia.org/wiki/Tilung_language}
\item Endangered Languages Project record: \url{https://www.endangeredlanguages.com/lang/4834}
\item Grambank entry (cites Opgenort's grammatical sketch): \url{https://grambank.clld.org/languages/tilu1238}
\item Glottolog: \url{https://glottolog.org/resource/languoid/id/tilu1238.bigmap.html}
\item Numerals page, Tilung: \url{https://lingweb.eva.mpg.de/channumerals/Tilung.htm}
\end{itemize}
The search summaries name these academic works (their own URLs were not seen): Opgenort (2011) on Tilung's position within Kiranti, in Himalayan Linguistics 10(1); Opgenort (2013), initial grammatical sketch; Rai (2012), phonological study, Nepalese Linguistics 27.

\section{Religious and Community Resources}
None found.

\section{Digital and Audio Resources}
None found.

\section{Curriculum Development Context}
With so few speakers, the work is documentation and revitalisation rather than a full K-12 curriculum. The academic sketches could underpin an orthography and word lists.

\end{document}